\documentclass{article}
\usepackage[utf8]{inputenc}
\usepackage[T1]{fontenc}
\usepackage{geometry}
\usepackage{graphicx}
\usepackage{float}
\usepackage{hyperref}
\usepackage[english,indonesian]{babel}

\title{Optimasi Pemrosesan Template LaTeX Menggunakan Pendekatan Berbasis AI}
\author{Academic Support Agent}
\date{\today}

\begin{document}
\maketitle

\begin{abstract}

\end{abstract}

\chapter{PENDAHULUAN}


\chapter{TINJAUAN PUSTAKA}


% Original content (commented out):
% \documentclass{article}
\usepackage{styles}

\title{Análise física sobre fênomenos de transporte devido à radiação eletromagnética}
\author{
    João \\
    Maria \\
    Alberto
}
\institutionname{Universidade Estadual de Campinas}
\coursename{F x29 - Física Experimental}


\begin{document}
% ---
% Capa e tabela de conteúdo
% ---
\printcover
{
    \hypersetup{allcolors=black}
    \makeatletter \renewcommand{\@dotsep}{4.5} \makeatother
    \tableofcontents
}

% ---
% Input dos capítulos
% ---
\newpage
\input{caps/root}

% ---
% Referências bibliográficas
% ---
\newpage
\printbibliography
\end{document}

\end{document}